\section{Adaptive Attack Resilience}
\label{sec:adaptive}

We evaluate defense resilience under adaptive attacks where the adversary has knowledge of the defense and iteratively optimizes payloads.
We use an attacker LLM (GPT-4o) that receives rejection feedback and refines injections over 10 campaigns $\times$ 5 iterations per defense.

\begin{table}[t]
\centering
\small
\caption{Adaptive attack resilience (GPT-4o attacker, 10 campaigns, 5 iterations each). Adaptive ASR = fraction of campaigns where the attacker eventually bypassed the defense. Resilience tier: \textbf{H}igh (0\%), \textbf{M}edium ($<$50\%), \textbf{L}ow ($\geq$50\%).}
\label{tab:adaptive}
\resizebox{\columnwidth}{!}{%
\begin{tabular}{l rr c}
\toprule
Defense & Adaptive ASR & Avg Iters & Tier \\
\midrule
D0 Baseline          & 0.300 & 3.8 & M \\
\midrule
D1 Spotlighting      & \textbf{0.000} & --- & \textbf{H} \\
D2 Policy Gate       & \textbf{0.000} & --- & \textbf{H} \\
D5 Sandwich          & \textbf{0.000} & --- & \textbf{H} \\
D6 Output Filter     & \textbf{0.000} & --- & \textbf{H} \\
D8 Semantic Firewall & \textbf{0.000} & --- & \textbf{H} \\
D10 \textbf{CIV}     & \textbf{0.000} & --- & \textbf{H} \\
\midrule
D7 Input Classifier  & 0.100 & 4.6 & M \\
D9 Dual-LLM          & 0.100 & 4.6 & M \\
D4 Re-execution      & 0.200 & 4.2 & M \\
D3 Task Alignment    & 0.400 & 3.4 & M \\
\bottomrule
\end{tabular}%
}
\begin{flushleft}
\scriptsize Defenses sorted by resilience. Six defenses achieve 0\% adaptive bypass; D10 (CIV) is the only multi-signal defense in this tier.
\end{flushleft}
\end{table}


\paragraph{Finding 10: Six Defenses Achieve Full Resilience.}
Six defenses achieve 0\% adaptive bypass: D1, D2, D5, D6, D8, and \defense{D10} (\civ{}) (Table~\ref{tab:adaptive}).
D3 shows the highest bypass (40\%), followed by D4 (20\%) and D7/D9 (10\%).
\civ{}'s resilience is notable because it operates through dynamic information-flow analysis rather than static rules---its three orthogonal signals create a combinatorial challenge: bypassing provenance requires using only goal-present entities, bypassing compatibility requires goal-aligned tools, and bypassing plan deviation requires staying within the inferred plan, which is fundamentally incompatible with executing a forbidden action.

\paragraph{The Adaptive Baseline Paradox.}
D0 achieves only 30\% adaptive bypass---lower than its 43.2\% static \asr{}.
Without a defense, there is no rejection feedback to guide payload refinement, so adaptive campaigns degrade to single-shot attacks.
This reveals that our adaptive evaluation measures \textbf{resistance to feedback-guided optimization}, not absolute adversarial robustness.
